\documentclass[11pt]{article}

\usepackage[a4paper,margin=1in]{geometry}
\usepackage{amsmath,amssymb}
\usepackage{booktabs}
\usepackage{longtable}
\usepackage{array}
\usepackage{hyperref}
\usepackage{siunitx}

\hypersetup{
  colorlinks=true,
  linkcolor=blue,
  urlcolor=blue
}

\title{High-Order Explicit Electro-Activation Solver with TNNP Ionic Model}
\author{cardiacFoam: \texttt{highOrderElectroActivationFoamExplicit}}
\date{May 6, 2026}

\begin{document}
\maketitle

\section{Purpose}

This report documents the solver
\path{/home/pablo/cardiacFoam-pc/applications/solvers/highOrderElectroActivationFoamExplicit/}.
The solver is an explicit monodomain electro-activation solver. It is intended
for benchmark-type simulations such as the Niederer et al. three-dimensional
cardiac electrophysiology benchmark. The solver supports either standard
finite-volume diffusion and cell-centred ionic integration, or high-order LRE
reconstruction for the diffusion operator and for the volume integration of
the ionic source.

The output quantities used by the benchmark are the activation time field,
the activation time at point P8, and activation time samples along the
diagonal from P1 to P8.

\section{Governing Equation}

The solver advances the monodomain equation for the transmembrane potential
$V_m$,
\begin{equation}
  \frac{\partial V_m}{\partial t}
  =
  \frac{1}{\chi C_m}\nabla\cdot(\boldsymbol{\sigma}\nabla V_m)
  - I_{\mathrm{ion}}(V_m,\mathbf{y})
  + \frac{I_{\mathrm{stim}}}{\chi C_m},
  \label{eq:monodomain}
\end{equation}
where $\chi$ is the membrane surface-to-volume ratio, $C_m$ is the membrane
capacitance per unit area, $\boldsymbol{\sigma}$ is the conductivity tensor,
$I_{\mathrm{stim}}$ is the applied stimulus current density, and
$\mathbf{y}$ denotes the ionic state vector.

The TNNP model uses membrane voltage in millivolts internally. The solver field
$V_m$ is passed to the ionic model and converted as
\begin{equation}
  V = 1000\,V_m,
\end{equation}
where $V$ is the voltage used in the TNNP equations below.

\section{Explicit Time Integration}

Let $K_i$ be a control volume. The explicit update used by the solver can be
written as
\begin{equation}
  V_{m,i}^{n+1}
  =
  V_{m,i}^{n}
  +
  \Delta t
  \left[
    \frac{1}{\chi C_m}L_i(V_m^n)
    - \overline{I}_{\mathrm{ion},i}^{\,n}
    + \frac{I_{\mathrm{stim},i}^{n}}{\chi C_m}
  \right],
  \label{eq:explicit-vm}
\end{equation}
where $L_i(V_m^n)$ is either the standard finite-volume diffusion operator or
the high-order LRE diffusion operator. The ionic state equations are advanced
after the voltage update through the selected OpenFOAM ODE solver. For the
Niederer setup used here, the selected ODE method is RKF45.

\section{High-Order Spatial Discretization}

\subsection{Diffusion Operator}

When \texttt{useHighOrder\_Vm false}, the diffusion term is evaluated with the
standard finite-volume operator
\begin{equation}
  L_i(V_m) = \left[\nabla\cdot(\boldsymbol{\sigma}\nabla V_m)\right]_i.
\end{equation}

When \texttt{useHighOrder\_Vm true}, LRE reconstruction is used at face
quadrature points. For a face $f$, the high-order diffusive flux is assembled
as
\begin{equation}
  \Phi_f
  =
  |S_f|
  \sum_{q=1}^{N_f}
  w_{fq}\,
  \mathbf{n}_f\cdot
  \left(
    \boldsymbol{\sigma}_{\mathrm{owner}(f)}
    \nabla_h V_m(\mathbf{x}_{fq})
  \right),
  \label{eq:ho-face-flux}
\end{equation}
and the cell diffusion contribution is
\begin{equation}
  L_i(V_m) =
  \frac{1}{|K_i|}
  \sum_{f\in\partial K_i}
  s_{if}\,\Phi_f,
\end{equation}
where $s_{if}=+1$ for outward owner-cell flux and $s_{if}=-1$ for neighbour
contributions.

\subsection{Ionic Source Quadrature}

The solver distinguishes the high-order reconstruction of the voltage from the
high-order volume integration of the ionic source.

If \texttt{useHighOrder\_Iion false}, the ionic model is integrated at cell
centres:
\begin{equation}
  \overline{I}_{\mathrm{ion},i}
  =
  I_{\mathrm{ion}}(V_{m,i},\mathbf{y}_i).
\end{equation}

If \texttt{useHighOrder\_Iion true}, the ionic states are stored and evolved at
the cell quadrature points $\mathbf{x}_{iq}$. The volume-averaged source is
then
\begin{equation}
  \overline{I}_{\mathrm{ion},i}
  =
  \frac{\sum_{q=1}^{N_i} w_{iq}
    I_{\mathrm{ion}}\left(V_{m,iq},\mathbf{y}_{iq}\right)}
       {\sum_{q=1}^{N_i} w_{iq}}.
  \label{eq:iion-quadrature}
\end{equation}
The quadrature-point voltage is reconstructed as
\begin{equation}
  V_{m,iq} =
  \begin{cases}
    R_{\mathrm{LRE}}[V_m](\mathbf{x}_{iq}),
      & \texttt{useHighOrder\_Vm true},\\
    V_{m,i} + \nabla V_{m,i}\cdot(\mathbf{x}_{iq}-\mathbf{x}_i),
      & \texttt{useHighOrder\_Vm false}.
  \end{cases}
\end{equation}
The ionic state vector at each quadrature point is advanced by
\begin{equation}
  \frac{d\mathbf{y}_{iq}}{dt}
  =
  \mathbf{F}_{\mathrm{TNNP}}\left(V_{iq},\mathbf{y}_{iq},t\right).
\end{equation}
For writing fields such as \texttt{Ca\_i}, the solver volume-averages the
quadrature states back to cells:
\begin{equation}
  \mathbf{y}_i
  =
  \frac{\sum_q w_{iq}\mathbf{y}_{iq}}{\sum_q w_{iq}}.
\end{equation}

\section{Boundary Conditions}

For the standard finite-volume branch, OpenFOAM boundary fields determine the
boundary contribution. For the high-order diffusion branch, \texttt{zeroGradient}
and \texttt{empty} patches are treated as true no-flux boundaries:
\begin{equation}
  \mathbf{n}\cdot\boldsymbol{\sigma}\nabla V_m = 0
  \quad\text{on}\quad \Gamma_N.
\end{equation}
In the implementation, this means that the high-order boundary face flux is set
to zero on these patches:
\begin{equation}
  \Phi_f = 0,
  \qquad f\subset\Gamma_N.
\end{equation}
This is important because simply reconstructing a high-order gradient at a
\texttt{zeroGradient} boundary can produce a non-zero normal flux unless the
Neumann condition is explicitly enforced at the high-order flux level.

Only \texttt{fixedValue} or \texttt{fixedVoltage} patches are included in the
high-order voltage stencil through boundary values. \texttt{zeroGradient}
patches are not included as value constraints in the high-order voltage stencil.

\section{Activation Time}

The activation time field is the first time at which $V_m$ crosses the
prescribed threshold $V_{\mathrm{thr}}$ from below. If
$V_{m,i}^n<V_{\mathrm{thr}}\le V_{m,i}^{n+1}$, the crossing time is linearly
interpolated as
\begin{equation}
  t_{\mathrm{act},i}
  =
  t^n
  +
  \Delta t
  \frac{V_{\mathrm{thr}}-V_{m,i}^n}
       {V_{m,i}^{n+1}-V_{m,i}^n}.
\end{equation}
The Niederer benchmark output point P8 is sampled from the activation-time
field at
\begin{equation}
  \mathbf{x}_{P8} = (0.02,\;0.003,\;0)\;{\rm m}.
\end{equation}

\section{TNNP Ionic Model}

The solver uses the ten Tusscher--Noble--Noble--Panfilov 2004 ionic model
through \texttt{ionicModel TNNP}. The state vector is
\begin{equation}
  \mathbf{y} =
  \left[
  V, K_i, Na_i, Ca_i, X_{r1}, X_{r2}, X_s, m, h, j,
  d, f, f_{Ca}, s, r, Ca_{SR}, g
  \right]^T .
\end{equation}
For the monodomain solver, $V$ is imposed by the PDE and is not evolved inside
the ionic ODE system unless \texttt{solveVmWithinODESolver} is enabled.

\subsection{Reversal Potentials}

\begin{align}
  E_{Na} &= \frac{RT}{F}\ln\left(\frac{Na_o}{Na_i}\right),\\
  E_K &= \frac{RT}{F}\ln\left(\frac{K_o}{K_i}\right),\\
  E_{Ks} &=
  \frac{RT}{F}
  \ln\left(
    \frac{K_o + P_{kna}Na_o}{K_i + P_{kna}Na_i}
  \right),\\
  E_{Ca} &= \frac{1}{2}\frac{RT}{F}\ln\left(\frac{Ca_o}{Ca_i}\right).
\end{align}

\subsection{Ionic Currents}

The total ionic current used by the model is
\begin{equation}
  I_{\mathrm{ion}} =
  i_{K1}+i_{to}+i_{Kr}+i_{Ks}+i_{CaL}+i_{NaK}
  +i_{Na}+i_{bNa}+i_{NaCa}+i_{bCa}+i_{pK}+i_{pCa}.
\end{equation}
The individual currents are
\begin{align}
  i_{Na} &= g_{Na}m^3hj(V-E_{Na}),\\
  i_{bNa} &= g_{bNa}(V-E_{Na}),\\
  i_{K1} &=
  g_{K1}x_{K1,\infty}\sqrt{\frac{K_o}{5.4}}(V-E_K),\\
  i_{to} &= g_{to}rs(V-E_K),\\
  i_{Kr} &= g_{Kr}\sqrt{\frac{K_o}{5.4}}X_{r1}X_{r2}(V-E_K),\\
  i_{Ks} &= g_{Ks}X_s^2(V-E_{Ks}),\\
  i_{bCa} &= g_{bCa}(V-E_{Ca}),\\
  i_{pK} &= \frac{g_{pK}(V-E_K)}{1+\exp((25-V)/5.98)},\\
  i_{pCa} &= \frac{g_{pCa}Ca_i}{Ca_i+K_{pCa}}.
\end{align}
The L-type calcium current is
\begin{equation}
  i_{CaL}
  =
  \frac{
  g_{CaL}df f_{Ca}
  \frac{4VF^2}{RT}
  \left(Ca_i\exp(2VF/RT)-0.341Ca_o\right)}
  {\exp(2VF/RT)-1}.
  \label{eq:ical}
\end{equation}
The sodium--potassium pump and sodium--calcium exchanger are
\begin{align}
  i_{NaK} &=
  \frac{
    P_{NaK}\frac{K_o}{K_o+K_{mk}}\frac{Na_i}{Na_i+K_{mNa}}
  }
  {
    1+0.1245\exp(-0.1VF/RT)+0.0353\exp(-VF/RT)
  },\\
  i_{NaCa} &=
  \frac{
    K_{NaCa}
    \left[
      \exp(\gamma V F/RT)Na_i^3 Ca_o
      -
      \exp((\gamma-1)VF/RT)Na_o^3 Ca_i\alpha
    \right]
  }
  {
    (K_{mNai}^3+Na_o^3)(K_{mCa}+Ca_o)
    \left[1+K_{sat}\exp((\gamma-1)VF/RT)\right]
  }.
\end{align}

\subsection{Gate Kinetics}

Most gates use the standard form
\begin{equation}
  \frac{dx}{dt} = \frac{x_\infty(V,Ca_i)-x}{\tau_x(V,Ca_i)}.
\end{equation}
The implemented steady states and time constants are
\begin{align}
  m_\infty &= \left[1+\exp((-56.86-V)/9.03)\right]^{-2},\\
  \alpha_m &= \left[1+\exp((-60-V)/5)\right]^{-1},\\
  \beta_m &=
  \frac{0.1}{1+\exp((V+35)/5)}
  +
  \frac{0.1}{1+\exp((V-50)/200)},\\
  \tau_m &= \alpha_m\beta_m,\\
  h_\infty &= \left[1+\exp((V+71.55)/7.43)\right]^{-2},\\
  \alpha_h &=
  \begin{cases}
  0.057\exp(-(V+80)/6.8), & V<-40,\\
  0, & V\ge -40,
  \end{cases}\\
  \beta_h &=
  \begin{cases}
  2.7\exp(0.079V)+310000\exp(0.3485V), & V<-40,\\
  \dfrac{0.77}{0.13[1+\exp((V+10.66)/-11.1)]}, & V\ge -40,
  \end{cases}\\
  \tau_h &= \frac{1}{\alpha_h+\beta_h},\\
  j_\infty &= h_\infty,\\
  \alpha_j &=
  \begin{cases}
  \dfrac{[-25428\exp(0.2444V)-6.948\cdot 10^{-6}\exp(-0.04391V)](V+37.78)}
  {1+\exp[0.311(V+79.23)]}, & V<-40,\\
  0, & V\ge -40,
  \end{cases}\\
  \beta_j &=
  \begin{cases}
  \dfrac{0.02424\exp(-0.01052V)}
  {1+\exp[-0.1378(V+40.14)]}, & V<-40,\\
  \dfrac{0.6\exp(0.057V)}
  {1+\exp[-0.1(V+32)]}, & V\ge -40,
  \end{cases}\\
  \tau_j &= \frac{1}{\alpha_j+\beta_j},\\
  d_\infty &= \left[1+\exp((-5-V)/7.5)\right]^{-1},\\
  \alpha_d &= \frac{1.4}{1+\exp((-35-V)/13)}+0.25,\\
  \gamma_{fCa} &= \frac{1.4}{1+\exp((V+5)/5)},\\
  \gamma_d &= \frac{1}{1+\exp((50-V)/20)},\\
  \tau_d &= \alpha_d\gamma_{fCa}+\gamma_d,\\
  f_\infty &= \left[1+\exp((V+20)/7)\right]^{-1},\\
  \tau_f &= 1125\exp[-(V+27)^2/240]+80
  +\frac{165}{1+\exp((25-V)/10)},\\
  X_{r1,\infty} &= \left[1+\exp((-26-V)/7)\right]^{-1},\\
  \alpha_{xr1} &= \frac{450}{1+\exp((-45-V)/10)},\\
  \beta_{xr1} &= \frac{6}{1+\exp((V+30)/11.5)},\\
  \tau_{xr1} &= \alpha_{xr1}\beta_{xr1},\\
  X_{r2,\infty} &= \left[1+\exp((V+88)/24)\right]^{-1},\\
  \alpha_{xr2} &= \frac{3}{1+\exp((-60-V)/20)},\\
  \beta_{xr2} &= \frac{1.12}{1+\exp((V-60)/20)},\\
  \tau_{xr2} &= \alpha_{xr2}\beta_{xr2},\\
  X_{s,\infty} &= \left[1+\exp((-5-V)/14)\right]^{-1},\\
  \alpha_{xs} &= 1100\left[1+\exp((-10-V)/6)\right]^{-1/2},\\
  \beta_{xs} &= \left[1+\exp((V-60)/20)\right]^{-1},\\
  \tau_{xs} &= \alpha_{xs}\beta_{xs},\\
  r_\infty &= \left[1+\exp((20-V)/6)\right]^{-1},\\
  \tau_r &= 9.5\exp[-(V+40)^2/1800]+0.8.
\end{align}
The $s$ gate depends on tissue type:
\begin{equation}
  s_\infty =
  \begin{cases}
  \left[1+\exp((V+20)/5)\right]^{-1},
  & \text{endocardial},\\
  1.1\left[1+\exp((V+28)/6)\right]^{-1},
  & \text{M cell or epicardial}.
  \end{cases}
\end{equation}
Its time constant is
\begin{equation}
  \tau_s =
  \begin{cases}
  85\exp[-(V+45)^2/320]
  +\dfrac{5}{1+\exp((V-20)/5)}+3,
  & \text{endocardial},\\
  1000\exp[-(V+67)^2/1000]+8,
  & \text{M cell or epicardial}.
  \end{cases}
\end{equation}
The inward rectifier auxiliary variable used by $i_{K1}$ is
\begin{align}
  \alpha_{K1}
  &= \frac{0.1}{1+\exp[0.06((V-E_K)-200)]},\\
  \beta_{K1}
  &=
  \frac{
  3\exp[0.0002((V-E_K)+100)]
  + \exp[0.1((V-E_K)-10)]}
  {1+\exp[-0.5(V-E_K)]},\\
  x_{K1,\infty}
  &= \frac{\alpha_{K1}}{\alpha_{K1}+\beta_{K1}}.
\end{align}
The calcium-dependent gates are
\begin{align}
  \alpha_{fCa} &= \left[1+\left(Ca_i/0.000325\right)^8\right]^{-1},\\
  \beta_{fCa} &= 0.1\left[1+\exp((Ca_i-0.0005)/0.0001)\right]^{-1},\\
  \beta_d &= 0.2\left[1+\exp((Ca_i-0.00075)/0.0008)\right]^{-1},\\
  f_{Ca,\infty} &= \frac{\alpha_{fCa}+\beta_{fCa}+\beta_d+0.23}{1.46}.
\end{align}
The $f_{Ca}$ rate is conditionally limited:
\begin{equation}
  \frac{df_{Ca}}{dt}
  =
  \begin{cases}
  0, & f_{Ca,\infty}>f_{Ca}\ \text{and}\ V>-60,\\
  (f_{Ca,\infty}-f_{Ca})/\tau_{fCa}, & \text{otherwise}.
  \end{cases}
\end{equation}
The RyR gate $g$ follows
\begin{align}
  g_\infty &=
  \begin{cases}
  \left[1+(Ca_i/0.00035)^6\right]^{-1}, & Ca_i<0.00035,\\
  \left[1+(Ca_i/0.00035)^{16}\right]^{-1}, & Ca_i\ge 0.00035,
  \end{cases}\\
  \frac{dg}{dt}
  &=
  \begin{cases}
  0, & g_\infty>g\ \text{and}\ V>-60,\\
  (g_\infty-g)/\tau_g, & \text{otherwise}.
  \end{cases}
\end{align}

\subsection{Concentration Equations}

The sodium and potassium concentration rates are
\begin{align}
  \frac{dNa_i}{dt}
  &=
  -\frac{(i_{Na}+i_{bNa}+3i_{NaK}+3i_{NaCa})C_m}{V_cF},\\
  \frac{dK_i}{dt}
  &=
  -\frac{[(i_{K1}+i_{to}+i_{Kr}+i_{Ks}+i_{pK}+I_{stim})
  -2i_{NaK}]C_m}{V_cF}.
\end{align}
The calcium release, uptake and leak terms are
\begin{align}
  i_{rel} &=
  \left[
    \frac{a_{rel}Ca_{SR}^2}{b_{rel}^2+Ca_{SR}^2}
    + c_{rel}
  \right]dg,\\
  i_{up} &= \frac{V_{max,up}}{1+K_{up}^2/Ca_i^2},\\
  i_{leak} &= V_{leak}(Ca_{SR}-Ca_i).
\end{align}
The buffered calcium rates are
\begin{align}
  J_{Ca_i} &=
  -\frac{(i_{CaL}+i_{bCa}+i_{pCa}-2i_{NaCa})C_m}
  {2V_cF}
  + i_{leak} - i_{up} + i_{rel},\\
  f_{JCa_i} &=
  \left[
    1+\frac{Buf_cK_{buf,c}}{(Ca_i+K_{buf,c})^2}
  \right]^{-1},\\
  \frac{dCa_i}{dt} &= J_{Ca_i}f_{JCa_i},\\
  J_{Ca_{SR}} &= \frac{V_c}{V_{SR}}(i_{up}-i_{rel}-i_{leak}),\\
  f_{JCa_{SR}} &=
  \left[
    1+\frac{Buf_{SR}K_{buf,SR}}{(Ca_{SR}+K_{buf,SR})^2}
  \right]^{-1},\\
  \frac{dCa_{SR}}{dt} &= J_{Ca_{SR}}f_{JCa_{SR}}.
\end{align}

\section{Numerical Protections in TNNP}

The generated TNNP algebra contains logarithms, powers and divisions that are
only defined for physically admissible states. The most sensitive quantities
are $Ca_i$, $Ca_{SR}$, $Na_i$, $K_i$, gate variables and the removable
singularity in Eq.~\eqref{eq:ical} near $V=0$.

The implemented protection is optional and controlled from
\texttt{cardiacProperties}. When \texttt{TNNPNumericalProtection true}, the
wrapper checks the state vector before calls to the generated TNNP algebra:
\begin{itemize}
  \item non-finite concentrations or concentrations below
  \texttt{TNNPConcentrationFloor} are replaced by that floor;
  \item gate variables are clipped to $[0,1]$ when
  \texttt{TNNPClampGates true};
  \item non-finite voltage is replaced by a resting fallback value;
  \item voltage is limited to
  $[V_{\min},V_{\max}]$ when \texttt{TNNPClampVmForRates true};
  \item if $|V|<\varepsilon_V$, $V$ is shifted to
  $\pm\varepsilon_V$ to avoid the direct evaluation of the
  $i_{CaL}$ denominator at zero.
\end{itemize}

Each correction is written to \texttt{TNNP\_numericalProtection.log} when
\texttt{TNNPLogNumericalProtection true}. The log columns are
\begin{equation}
  t_{\mathrm{ms}},\quad t_{\mathrm{s}},\quad
  \text{integration point},\quad
  \text{phase},\quad
  \text{variable},\quad
  \text{old value},\quad
  \text{corrected value},\quad
  \text{reason}.
\end{equation}
This makes the protection auditable: the run can be continued without losing
information about which state became non-physical and where the correction was
applied.

\section{Niederer P8 Activation-Time Results}
\label{sec:results}

This section reports two-dimensional Niederer-benchmark results read from
both result trees of this solver:
\path{/home/pablo/OpenFOAM/pablo-v2312/run/tutorials_electro/highOrderNiedererEtAl2012Explicit/}
and
\path{/home/pablo/OpenFOAM/pablo-v2312/run/tutorials_electro/highOrderNiedererEtAl2012Explicit_2/}.
The first folder contains the historical sweep over the requested
$\Delta t \in \{10^{-2}, 5\!\cdot\!10^{-5}, 10^{-5}, 5\!\cdot\!10^{-6}\}$~s;
the second folder is a more recent and more complete sweep at the
single requested $\Delta t = 10^{-2}$~s. Because the explicit solver
applies the CFL stability limit, the value the solver actually uses
($\Delta t_{\rm eff}$, table at the top of the next subsection) is far
smaller than the requested one and depends on the mesh size~$\Delta x$.
Computational cost data come from \texttt{/usr/bin/time -v} and are
reported as elapsed wall-clock time and peak resident set size (RSS).
\textsc{NC} means that no run was found for that combination in the
results tree.

\noindent The label \emph{Vm} refers to the LRE p-level used for the
high-order diffusion reconstruction, and the label \emph{Iion} refers
to the LRE p-level used for the cell-quadrature ionic-source
integration. \texttt{NO} means that high-order is disabled for that
quantity (cell-centre Laplacian / cell-centre ionic source).

% Auto-generated by make_latex_tables.py
\subsection{P8 activation time at the requested $\Delta t = 1.0\cdot10^{-2}$~s}
Effective time step set by the explicit CFL limit:
\begin{center}\small
\begin{tabular}{cccccc}\toprule
$\Delta x$ [mm] & 0.5 & 0.25 & 0.125 & 0.1 & 0.05\\
$\Delta t_{\rm eff}$ [s] & $1.3\!\cdot\!10^{-4}$ & $3.3\!\cdot\!10^{-5}$ & $8.2\!\cdot\!10^{-6}$ & $5.2\!\cdot\!10^{-6}$ & $1.3\!\cdot\!10^{-6}$\\
\bottomrule\end{tabular}\end{center}

\begin{center}\footnotesize
\begin{longtable}{llrrrrr}
\caption{P8 activation time [ms]: rows are LRE p-levels for Vm $\times$ Iion; columns are mesh sizes $\Delta x$ (mm). \textsc{NC} = no run data found.}\\
\toprule
Vm & Iion & $0.5$ & $0.25$ & $0.125$ & $0.1$ & $0.05$\\
\midrule\endfirsthead\toprule
Vm & Iion & $0.5$ & $0.25$ & $0.125$ & $0.1$ & $0.05$\\
\midrule\endhead
NO & NO & 160.63 & 68.12 & 50.35 & 48.15 & 45.62\\
NO & p1 & 96.43 & 59.69 & 48.33 & 47.08 & 45.53\\
NO & p2 & 34.71 & 42.82 & 44.70 & 44.95 & NC\\
NO & p3 & 35.33 & 43.46 & 44.67 & 44.92 & NC\\\hline
p1 & NO & 14.25 & 75.16 & 54.23 & 51.20 & 46.59\\
p1 & p1 & 104.08 & 65.69 & 51.94 & 49.91 & 46.47\\
p1 & p2 & 38.45 & 47.47 & 47.24 & 46.97 & 46.02\\
p1 & p3 & 40.65 & 48.30 & 47.22 & 46.94 & NC\\\hline
p2 & NO & 0.00 & 89.11 & 60.13 & 55.60 & 47.85\\
p2 & p1 & 159.22 & 74.08 & 56.36 & 53.33 & 47.58\\
p2 & p2 & 42.88 & 48.89 & 48.50 & 48.14 & 46.56\\
p2 & p3 & 48.72 & 50.35 & 48.65 & 48.19 & NC\\\hline
p3 & NO & 192.55 & 73.84 & 54.51 & 51.31 & 46.24\\
p3 & p1 & 67.83 & 62.86 & 51.48 & 49.52 & 46.07\\
p3 & p2 & 34.76 & 43.10 & 45.51 & 45.70 & NC\\
p3 & p3 & 36.02 & 43.80 & 45.54 & 45.69 & NC\\
\bottomrule\end{longtable}\end{center}

\subsection{Wall-clock cost per configuration (minutes)}
Total solver elapsed time as reported by \texttt{/usr/bin/time -v}.
\begin{center}\footnotesize
\begin{longtable}{llrrrrr}
\caption{Solver wall-clock time [min].}\\
\toprule
Vm & Iion & $0.5$ & $0.25$ & $0.125$ & $0.1$ & $0.05$\\
\midrule\endfirsthead\toprule
Vm & Iion & $0.5$ & $0.25$ & $0.125$ & $0.1$ & $0.05$\\
\midrule\endhead
NO & NO & 1 & 2 & 9 & 14 & 168\\
NO & p1 & 1 & 3 & 15 & 25 & 416\\
NO & p2 & 1 & 6 & 42 & 73 & NC\\
NO & p3 & 1 & 9 & 59 & 100 & NC\\\hline
p1 & NO & 1 & 2 & 9 & 15 & 203\\
p1 & p1 & 1 & 3 & 17 & 27 & 381\\
p1 & p2 & 1 & 7 & 45 & 77 & 1059\\
p1 & p3 & 2 & 9 & 60 & 108 & NC\\\hline
p2 & NO & 1 & 2 & 10 & 16 & 239\\
p2 & p1 & 1 & 4 & 19 & 31 & 398\\
p2 & p2 & 1 & 7 & 47 & 77 & 1167\\
p2 & p3 & 2 & 10 & 62 & 113 & NC\\\hline
p3 & NO & 1 & 2 & 9 & 15 & 197\\
p3 & p1 & 1 & 3 & 18 & 31 & 402\\
p3 & p2 & 1 & 6 & 47 & 75 & NC\\
p3 & p3 & 1 & 9 & 58 & 106 & NC\\
\bottomrule\end{longtable}\end{center}

\subsection{Peak resident set size per configuration (MiB)}
Maximum physical RAM used by the solver process.
\begin{center}\footnotesize
\begin{longtable}{llrrrrr}
\caption{Peak RSS [MiB].}\\
\toprule
Vm & Iion & $0.5$ & $0.25$ & $0.125$ & $0.1$ & $0.05$\\
\midrule\endfirsthead\toprule
Vm & Iion & $0.5$ & $0.25$ & $0.125$ & $0.1$ & $0.05$\\
\midrule\endhead
NO & NO & 148 & 152 & 165 & 179 & 297\\
NO & p1 & 149 & 155 & 182 & 202 & 388\\
NO & p2 & 152 & 168 & 234 & 287 & NC\\
NO & p3 & 154 & 174 & 261 & 329 & NC\\\hline
p1 & NO & 148 & 154 & 177 & 198 & 377\\
p1 & p1 & 150 & 158 & 200 & 234 & 516\\
p1 & p2 & 153 & 170 & 254 & 318 & 855\\
p1 & p3 & 155 & 176 & 282 & 361 & NC\\\hline
p2 & NO & 149 & 155 & 180 & 204 & 398\\
p2 & p1 & 150 & 159 & 217 & 260 & 622\\
p2 & p2 & 153 & 173 & 271 & 345 & 961\\
p2 & p3 & 155 & 180 & 298 & 388 & NC\\\hline
p3 & NO & 150 & 160 & 202 & 236 & 524\\
p3 & p1 & 151 & 174 & 277 & 354 & 995\\
p3 & p2 & 155 & 188 & 331 & 439 & NC\\
p3 & p3 & 156 & 195 & 356 & 482 & NC\\
\bottomrule\end{longtable}\end{center}

\section{Implementation Notes}

The most relevant implementation details are:
\begin{itemize}
  \item \texttt{useHighOrder\_Vm} controls the LRE high-order diffusion
  reconstruction.
  \item \texttt{useHighOrder\_Iion} controls whether TNNP states are evolved
  at cell quadrature points and whether $I_{\mathrm{ion}}$ is volume-integrated
  with those points.
  \item \texttt{useHighOrder\_states} is still accepted as a legacy dictionary
  keyword, but the current solver interprets it as the high-order
  $I_{\mathrm{ion}}$ quadrature switch.
  \item \texttt{activationTime} is written as a field and can be viewed in
  ParaView.
  \item \texttt{P8\_activationTime.dat} and
  \texttt{diagonalActivationTime.dat} are written for benchmark postprocessing.
\end{itemize}

\end{document}
